\documentclass[12pt]{article}
\usepackage[margin=1in]{geometry}
\usepackage{titlesec}
\usepackage{hyperref}
\usepackage{booktabs}
\usepackage{array}
\usepackage{longtable}
\usepackage{parskip}

\title{Software Requirements Specification\\
\large QTC Workflow Queue (WFQ)}
\author{Original Authors: Dylan Huang, Laila Velasquez, Samuel Acevedo, Haik Oganesyan,\\
Ruben Flores, Slok Patel, Edwin Rojas, Jesus Ojeda, Alexander Diaz, Zhen Zhao\\[0.5em]
\textit{Adapted to \LaTeX{} by CS3338 Course Team}}
\date{Version 1.3 \\ Originally: October 17, 2025 \\ Adapted: \today}

\begin{document}

\maketitle
\newpage
\tableofcontents
\newpage

% ─────────────────────────────────────────────
\section*{Revision History}
% ─────────────────────────────────────────────
\begin{center}
\begin{tabular}{llll}
\toprule
Name & Date & Reason For Changes & Version \\
\midrule
All group members  & 10/17/2025 & Initial Document                          & 1.0 \\
Slok Patel         & 10/17/2025 & Section 1 and 2 introduction              & 1.1 \\
Zhen Zhao          & 10/22/2025 & Section 4 Requirement Specifications      & 1.2 \\
Edwin Rojas        & 10/23/2025 & Section 3 External Interface Requirements & 1.2 \\
Slok Patel         & 10/23/2025 & Section 4, 5, \& Appendix (ABC)          & 1.2 \\
Dylan Huang        & 10/23/2025 & Section 6 Legal and Ethical Considerations & 1.2 \\
Laila Velasquez    & 10/23/2025 & Table of Contents                         & 1.2 \\
Edwin Rojas        & 11/4/2025  & Section 1, 2, and 3 Revisions            & 1.3 \\
\bottomrule
\end{tabular}
\end{center}
\newpage

% ─────────────────────────────────────────────
\section{Introduction}
% ─────────────────────────────────────────────

\subsection{Purpose}
The purpose of this Software Requirements Specification (SRS) document is to define the
functional and nonfunctional requirements for the Workflow Queue (WFQ) Administrator
Application. This document specifies what the software will do and the conditions under which
it will operate. It serves as a guide for developers, testers, and stakeholders by clearly
outlining the system's intended features and behaviors.

The WFQ application will be developed using ASP.NET MVC with C\#, Entity Framework, and
stored procedures, with both SQL Server and Oracle serving as supported backends. This SRS
focuses on Version 1.0 of the new Administrator Application, which will be developed using
ASP.NET MVC with C\# and Entity Framework. Its sole backend interaction will be with the
existing QTC-provided SQL Server database. This document provides the foundation for the
design and implementation activities described in the Software Design Document (SDD).

\subsection{Intended Audience and Reading Suggestions}
This document is intended for a variety of audiences involved in the project lifecycle:

\begin{itemize}
    \item \textbf{Project Sponsor (Leidos QTC):} Can use this document to verify that the
    planned application meets their specified needs and to understand the project's scope and
    boundaries.
    \item \textbf{Developers:} Should focus on the detailed functional and non-functional
    requirements in Sections 4 and 5 to guide implementation.
    \item \textbf{QA Testers:} Should use the requirements in Section 4 as a basis for writing
    test cases and defining acceptance criteria.
    \item \textbf{Project Managers:} Can use the overall product description and defined scope
    to manage project execution and stakeholder expectations.
\end{itemize}

\subsection{Product Scope}
The WFQ Administrator Application is a new, web-based tool designed for authorized technical
administrators to create, configure, and manage the workflows used by the main QTC WFQ
system. The application will achieve this by directly reading from and writing to the existing
WFQ SQL Server database.

The core purpose of this application is to provide a user-friendly graphical interface for
managing the underlying data that defines the workflow engine's behavior. The initial release
will focus on providing full administrative control over workflows and their associated steps.

\textbf{IN SCOPE:}
\begin{itemize}
    \item Designing, building, and delivering a new, standalone Administrator web application.
    \item Creating a data access layer within the new application that communicates directly
    with the existing WFQ SQL Server database.
    \item Implementing all administrative features as specified in the ``Workflow Queue
    Administrator Requirements'' document.
    \item Implementing role-based access control within the new application based on a URL
    query string parameter.
\end{itemize}

\textbf{OUT OF SCOPE:}
\begin{itemize}
    \item Accessing, modifying, or using any source code from the existing live WFQ
    application.
    \item Calling, updating, or interacting with any pre-existing backend REST API.
    \item Building or modifying the existing WFQ User \& Manager interface.
    \item Integrating with external third-party services or creating advanced analytics
    dashboards.
\end{itemize}

\subsection{Definitions, Acronyms, and Abbreviations}
\begin{itemize}
    \item \textbf{WFQ} -- Workflow Queue
    \item \textbf{ASP.NET MVC} -- A Model-View-Controller framework for web applications
    using Microsoft's .NET platform
    \item \textbf{C\#} -- Programming language used for the system's logic and business layer
    \item \textbf{SQL Server / Oracle} -- Database management systems used for backend storage
    \item \textbf{Entity Framework (EF)} -- An object-relational mapping (ORM) framework used
    for data operations
    \item \textbf{QA} -- Quality Assurance
    \item \textbf{PO} -- Product Owner
    \item \textbf{SRS} -- Software Requirements Specification
    \item \textbf{SDD} -- Software Design Document
\end{itemize}

\subsection{References}
\begin{itemize}
    \item Microsoft. ``C\# Coding Conventions.'' \url{https://learn.microsoft.com/en-us/dotnet/csharp/fundamentals/coding-style/coding-conventions}
    \item Microsoft Learn. ``Build Web APIs with ASP.NET Core.'' \url{https://learn.microsoft.com/en-us/training/modules/build-web-api-aspnet-core/}
    \item Microsoft Learn. ``Install SQL Server Management Studio (SSMS).'' \url{https://learn.microsoft.com/en-us/ssms/install/install}
    \item Oracle. ``SQL Developer Download.'' \url{https://www.oracle.com/database/sqldeveloper/technologies/download/}
\end{itemize}

% ─────────────────────────────────────────────
\section{Overall Description}
% ─────────────────────────────────────────────

The WFQ (Workflow Queue) application provides a centralized platform for automating and
managing multi-step workflows across organizational departments. Its purpose is to replace
inefficient manual processes and disconnected tools with a cohesive system that improves
coordination, reduces delays, and enhances visibility.

\subsection{System Analysis}
The WFQ system is designed to solve the problem of inefficient and fragmented workflow
management. Many organizations currently use manual methods or isolated applications for
approvals, which often result in miscommunication, data loss, and process delays. The WFQ
application automates these processes by enabling users to create, track, and complete
workflows through a unified interface.

The main goals of the system include increasing process transparency, reducing approval times,
and ensuring accountability at every step. Key technical challenges involve maintaining data
integrity between SQL Server and Oracle databases and implementing secure, role-based
authentication.

\subsection{Product Perspective}
The WFQ Administrator Application is a new, self-contained, internal web tool. It is not an
enhancement to the existing WFQ application but rather a separate, complementary system that
shares a common database.

Its position within the QTC ecosystem is that of a ``back-office'' or ``control panel'' tool.
While the primary WFQ application is used by front-line staff to execute work, this new
Administrator Application will be used by technical staff to define and manage the structure
of that work.

\subsection{Product Functions}
The primary functions of the WFQ Administrator Application include:
\begin{itemize}
    \item \textbf{Workflow Creation and Configuration:} Allows authorized administrators to
    define new workflows.
    \item \textbf{Workflow Step Management:} Allows administrators to add, edit, and delete
    the individual steps that constitute a workflow, including their sequence and properties.
    \item \textbf{Task Oversight:} Provides administrators with the ability to view tasks
    within a workflow and perform limited management functions, such as reassigning or
    completing a task.
    \item \textbf{History and Auditing:} Displays a history of tasks and changes within a
    workflow to support troubleshooting and oversight.
\end{itemize}

\subsection{User Classes and Characteristics}
The WFQ system will be used by multiple user groups with different roles and privileges:
\begin{itemize}
    \item \textbf{Administrators:} Manage user accounts, permissions, and system
    configurations. They are responsible for defining workflow templates and monitoring
    overall system health.
    \item \textbf{Managers / Approvers:} Review and approve workflow items. They have access
    to reporting tools and can delegate or escalate tasks when necessary.
    \item \textbf{Regular Users / Submitters:} Initiate workflows, upload relevant documents,
    and monitor the progress of their requests.
    \item \textbf{QA Testers:} Test application features and ensure compliance with
    functionality, performance, and security standards.
    \item \textbf{Developers:} Maintain and enhance the system, ensuring compatibility across
    environments.
\end{itemize}

\subsection{Operating Environment}
The WFQ application will operate in a web-based environment using the following components:
\begin{itemize}
    \item \textbf{Frontend:} Developed with ASP.NET MVC for the UI layer, compatible with
    modern browsers (Chrome, Edge, Firefox).
    \item \textbf{Backend:} Built in C\# with Entity Framework, handling data access logic.
    \item \textbf{Databases:} Microsoft SQL Server and Oracle Database for backend storage and
    stored procedure execution.
    \item \textbf{Operating System:} Hosted on Windows Server with support for Internet
    Information Services (IIS).
    \item \textbf{Development Tools:} Visual Studio, SQL Server Management Studio (SSMS), and
    Oracle SQL Developer.
\end{itemize}

\subsection{Design and Implementation Constraints}
\begin{itemize}
    \item \textbf{Direct Database Access Only:} The application must not attempt to interface
    with the existing WFQ application at the code or API level. All interactions must occur
    via direct queries and commands to the SQL Server database.
    \item \textbf{Technology Stack:} The project is constrained to the existing technology
    stack (C\#, ASP.NET, MS SQL Server) to ensure maintainability and consistency.
    \item \textbf{Security \& Compliance (HIPAA):} The application will be handling data that
    is subject to HIPAA regulations. All access must be strictly controlled, and all
    administrative actions that modify data must be auditable.
    \item \textbf{Authorization Model:} The application must rely solely on a URL query string
    parameter for determining user authorization. It will not integrate with a separate
    authentication service API.
\end{itemize}

\subsection{User Documentation}
Upon completion, the WFQ application will include:
\begin{itemize}
    \item \textbf{User Manual:} Provides step-by-step guidance for creating, managing, and
    approving workflows.
    \item \textbf{Online Help Guides:} Integrated help documentation within the application's
    UI.
    \item \textbf{Developer Documentation:} Includes database schema details and setup
    instructions.
    \item \textbf{Training Materials:} Optional tutorial videos or PDF guides.
\end{itemize}

\subsection{Assumptions and Dependencies}
\begin{itemize}
    \item \textbf{Stable Database Schema:} The schema of the existing WFQ database is stable
    and will not undergo significant changes during the project lifecycle.
    \item \textbf{Database Accessibility:} The development team will be provided with network
    access and credentials to a development/testing instance of the WFQ database.
    \item \textbf{Sponsor-Provided Authorization:} The hosting environment will provide a
    reliable and secure mechanism for passing the user's administrative status via the URL
    query string.
    \item \textbf{Team Skillset:} The team possesses, or can acquire, the necessary skills in
    C\#, ASP.NET MVC, and SQL to complete the project.
\end{itemize}

\subsection{Apportioning of Requirements}
Certain advanced features may be reserved for future versions, including:
\begin{itemize}
    \item Integration with third-party APIs for document signing.
    \item Advanced analytics dashboards using Power BI.
    \item Mobile version of the WFQ system for iOS and Android platforms.
    \item Support for multilingual UI and accessibility enhancements.
\end{itemize}

% ─────────────────────────────────────────────
% SNAPSHOT 2 ADDITION
% ─────────────────────────────────────────────
\subsection{Snapshot 2 Addition: Email Notification System}
\textit{Added by CS3338 Course Team --- Snapshot 2.}

The Email Notification System is a new feature added in Snapshot 2. It extends the existing
WFQ Administrator Application with automated email alerts that notify relevant users when
key workflow events occur. This feature was originally identified in Section 2.9 of this
document as a planned future addition and is now being formally specified.

\subsubsection{Purpose}
To improve workflow visibility and reduce manual follow-up by automatically notifying
assigned users and administrators when tasks are assigned, reassigned, or completed, and
when new workflow steps are created.

\subsubsection{Functional Requirements --- Notification Triggers}
\begin{itemize}
    \item \textbf{NFR-1:} The system shall send an email notification to the assigned user
    when a task is assigned or reassigned to them.
    \item \textbf{NFR-2:} The system shall send an email notification to the workflow
    administrator when a task is marked as completed.
    \item \textbf{NFR-3:} The system shall send an email notification to the workflow
    administrator when a new workflow step is created within a workflow they manage.
    \item \textbf{NFR-4:} The system shall not send a notification for read-only operations
    (e.g., viewing a workflow or step).
    \item \textbf{NFR-5:} If a notification fails to send (e.g., SMTP error), the system
    shall log the failure and continue normal operation without interrupting the user's
    workflow action.
\end{itemize}

\subsubsection{Functional Requirements --- Email Content}
\begin{itemize}
    \item \textbf{NFR-6:} Each notification email shall include: the event type, the
    workflow name, the step name (if applicable), the task details (if applicable), and a
    direct URL link to the relevant workflow in the application.
    \item \textbf{NFR-7:} Emails shall be sent from a configured sender address defined in
    the application's server-side configuration (\texttt{appsettings.json}).
    \item \textbf{NFR-8:} Email subjects shall clearly identify the event type (e.g.,
    ``[WFQ] Task Assigned: \{TaskName\}'').
\end{itemize}

\subsubsection{Functional Requirements --- Configuration}
\begin{itemize}
    \item \textbf{NFR-9:} SMTP host, port, sender address, and credentials shall be
    configurable via \texttt{appsettings.json} and shall not be hardcoded.
    \item \textbf{NFR-10:} SMTP credentials shall not be committed to source control.
    They shall be managed via environment variables or a secrets manager.
    \item \textbf{NFR-11:} The notification feature shall be toggleable via a configuration
    flag (\texttt{Notifications:Enabled: true/false}) to allow disabling in development or
    testing environments.
\end{itemize}

\subsubsection{Non-Functional Requirements}
\begin{itemize}
    \item \textbf{NFR-12:} Notification sending shall be performed asynchronously and must
    not block or delay the user's primary action (e.g., task assignment).
    \item \textbf{NFR-13:} The notification service shall be fully unit-testable via
    dependency injection and interface mocking.
\end{itemize}

\subsubsection{New Dependencies}
\begin{itemize}
    \item \textbf{MailKit} (NuGet, pinned version) --- cross-platform .NET SMTP library
    used for composing and sending email notifications.
\end{itemize}

% ─────────────────────────────────────────────
% SNAPSHOT 3 ADDITION
% ─────────────────────────────────────────────
\subsection{Snapshot 3 Addition: Workflow Export / Report Generation}
\textit{Added by CS3338 Course Team --- Snapshot 3.}

The Workflow Export and Report Generation feature allows administrators to download a
workflow's full details as a PDF or CSV file directly from the Workflow Details page. This
is a new feature not present in the original project.

\subsubsection{Purpose}
To give administrators a portable, shareable record of a workflow's current state ---
including its steps, tasks, and metadata --- without requiring direct database access or
manual data collection.

\subsubsection{Functional Requirements --- PDF Export}
\begin{itemize}
    \item \textbf{EXP-1:} The system shall allow an admin to export a workflow as a
    formatted PDF file from the Workflow Details page.
    \item \textbf{EXP-2:} The PDF shall include: workflow name, description, status,
    created date, a table of all steps (StepNumber, StepName, Description), and a table
    of all tasks (AssignedTo, Status, Step).
    \item \textbf{EXP-3:} The PDF file shall be named
    \texttt{\{WorkflowName\}\_Report\_\{Date\}.pdf}.
    \item \textbf{EXP-4:} The PDF shall be generated server-side and delivered as a
    file download response. The user shall not be navigated away from the Workflow Details
    page.
    \item \textbf{EXP-5:} The PDF shall be readable and correctly formatted on standard
    PDF viewers (Adobe Acrobat, browser PDF viewer).
\end{itemize}

\subsubsection{Functional Requirements --- CSV Export}
\begin{itemize}
    \item \textbf{EXP-6:} The system shall allow an admin to export a workflow's task list
    as a CSV file from the Workflow Details page.
    \item \textbf{EXP-7:} The CSV shall include the following columns:
    WorkflowName, StepName, StepNumber, TaskID, AssignedTo, Status, StartDate,
    CompletionDate.
    \item \textbf{EXP-8:} The CSV file shall be named
    \texttt{\{WorkflowName\}\_Tasks\_\{Date\}.csv}.
    \item \textbf{EXP-9:} The CSV shall use UTF-8 encoding and standard comma-delimited
    format compatible with Microsoft Excel and Google Sheets.
    \item \textbf{EXP-10:} The CSV shall be delivered as a file download response without
    navigating away from the Workflow Details page.
\end{itemize}

\subsubsection{Functional Requirements --- General}
\begin{itemize}
    \item \textbf{EXP-11:} Export actions shall require the same admin authorization as
    all other actions in the application (\texttt{[Authorize(Policy = "IsAdmin")]}).
    \item \textbf{EXP-12:} If a workflow has no steps or no tasks, the export shall still
    succeed and the relevant sections shall display ``No data available.''
    \item \textbf{EXP-13:} Export operations shall not modify any data in the database.
    They are strictly read-only operations.
\end{itemize}

\subsubsection{Non-Functional Requirements}
\begin{itemize}
    \item \textbf{EXP-14:} PDF and CSV generation shall complete within 3 seconds for
    workflows with up to 100 steps and 500 tasks under normal server load.
    \item \textbf{EXP-15:} The export service shall be fully unit-testable via dependency
    injection and interface mocking.
\end{itemize}

\subsubsection{New Dependencies}
\begin{itemize}
    \item \textbf{QuestPDF} (NuGet, pinned version) --- a modern .NET library for
    server-side PDF generation. Used for PDF export functionality.
    \item \textbf{No additional packages for CSV} --- CSV generation uses built-in
    \texttt{System.Text} and \texttt{StringBuilder} classes.
\end{itemize}

% ─────────────────────────────────────────────
\section{External Interface Requirements}
% ─────────────────────────────────────────────

\subsection{User Interfaces}
The system will present a single, new web-based user interface — the \textbf{Administrator UI}.
This is a new, dedicated web application built from scratch for technical administrators to
manage the workflow system. It will provide all functionality specified in the ``Workflow Queue
Administrator Requirements'' document, including the ability to create, edit, and view
workflows and their steps.

\textbf{GUI and Style Standards:}
\begin{itemize}
    \item \textbf{Branding:} All UIs must adhere to the Leidos corporate branding guidelines,
    including the official color palette and logos.
    \item \textbf{Layout:} Screen layouts will be standardized. Standard functions (e.g.,
    search, create, edit) will be presented consistently.
    \item \textbf{Error Display:} The system will display user-friendly error messages for
    invalid inputs or exceptions.
    \item \textbf{Accessibility:} The UI design should consider conformance with Section 508
    compliance rules.
\end{itemize}

\subsection{Hardware Interfaces}
The WFQ Administrator Application is a software-only system. It does not interface directly
with any custom or specialized hardware. All operations are handled through standard user
input devices (keyboard, mouse).

\subsection{Software Interfaces}
The system has one primary external software interface: the \textbf{Microsoft SQL Server
Database}. This is the existing, persistent data store for the live QTC WFQ system and is the
single source of truth for all workflow, step, and task data. Data access will be managed via
Entity Framework and by executing existing stored procedures where applicable.

\subsection{Communications Interfaces}
\begin{itemize}
    \item \textbf{Protocol:} All communication between the administrator's web browser and the
    application's web server (IIS) shall be conducted securely over HTTPS/TLS.
    \item \textbf{Authentication:} User authentication is assumed to be handled by the hosting
    environment prior to the user accessing the application's URL. The application itself does
    not manage user logins or passwords.
\end{itemize}

% ─────────────────────────────────────────────
\section{Requirements Specification}
% ─────────────────────────────────────────────

\subsection{Functional Requirements}

\subsubsection{Authentication and Authorization}
\begin{itemize}
    \item 4.1.1.1 The system shall require all users to authenticate before accessing the
    Workflow Queue Administrator interface.
    \item 4.1.1.2 The system shall verify the user's authorization level upon authentication.
    \item 4.1.1.3 The system shall restrict access to users with the ``Admin'' role.
    \item 4.1.1.4 The system shall deny access to users who do not have admin privileges,
    displaying an ``Access Denied'' message.
    \item 4.1.1.5 The system shall validate an admin authorization flag via the query string
    before granting access.
\end{itemize}

\subsubsection{Workflow Management}

\textbf{Workflow Display}
\begin{itemize}
    \item 4.2.1.1 The system shall display a paginated list of all existing workflows.
    \item 4.2.1.2 Each workflow entry shall display: Name, Description, Status (New, In
    Progress, Completed), and Actions (Edit, View Steps).
    \item 4.2.1.3 The system shall allow admins to search or filter workflows by Name.
    \item 4.2.1.4 The system shall allow sorting by Name and CreatedDate.
\end{itemize}

\textbf{Workflow Creation}
\begin{itemize}
    \item 4.2.2.1 The system shall allow an admin to create a new workflow.
    \item 4.2.2.2 The creation form shall include Name (required, unique) and Description
    (optional).
    \item 4.2.2.3 The system shall enforce a maximum length for the Name field.
    \item 4.2.2.4 The system shall validate that workflow names are unique.
\end{itemize}

\textbf{Workflow Update}
\begin{itemize}
    \item 4.2.3.1 The system shall allow admins to update existing workflows.
    \item 4.2.3.2 The system shall prepopulate the update form with current workflow data.
    \item 4.2.3.3 The same validation rules as creation shall apply.
    \item 4.2.3.4 The system shall not allow workflow deletion through the interface.
    \item 4.2.3.5 The system shall display an informational note that deletion must be handled
    manually if required.
\end{itemize}

\subsubsection{Workflow Step Management}

\textbf{Step Display}
\begin{itemize}
    \item 4.3.1.1 The system shall display all steps associated with a selected workflow.
    \item 4.3.1.2 Each step shall display: StepNumber, StepName, Description, and Actions
    (Edit/Delete).
\end{itemize}

\textbf{Step Creation}
\begin{itemize}
    \item 4.3.2.1 The system shall allow the admin to create new steps within a workflow.
    \item 4.3.2.2 Required fields: StepName (unique within workflow).
    \item 4.3.2.3 Optional fields: StepNumber (auto-assigned if omitted).
    \item 4.3.2.4 The system shall validate StepNumber uniqueness within the workflow.
\end{itemize}

\textbf{Step Update}
\begin{itemize}
    \item 4.3.3.1 The system shall allow admins to update existing step data.
    \item 4.3.3.2 The system shall prepopulate the step form with current information.
    \item 4.3.3.3 Validation rules shall match those used in creation.
\end{itemize}

\textbf{Step Deletion}
\begin{itemize}
    \item 4.3.4.1 The system shall prompt for confirmation before deleting a step.
    \item 4.3.4.2 The system shall check for dependent tasks before deletion.
    \item 4.3.4.3 If dependencies exist, the system shall reassign tasks to the next step
    before deletion.
    \item 4.3.4.4 The system shall delete steps using Entity Framework, preventing cascading
    deletions beyond the step.
\end{itemize}

\subsubsection{Task Management}
\begin{itemize}
    \item 4.4.1 The system shall display all tasks associated with a selected workflow.
    \item 4.4.2 Each task shall display: AssignedTo, Status, Step, and Actions (Complete Task
    / Reassign).
    \item 4.4.3 The system shall allow admins to assign or unassign tasks to users.
    \item 4.4.4 The system shall allow admins to mark tasks as complete.
\end{itemize}

\subsubsection{Error Handling and Validation}
\begin{itemize}
    \item 4.5.1 The system shall perform input validation on both client and server sides.
    \item 4.5.2 The system shall display user-friendly error messages for invalid inputs.
    \item 4.5.3 The system shall handle and report database exceptions such as concurrency or
    foreign key violations.
    \item 4.5.4 Invalid or incomplete form submissions shall be rejected.
\end{itemize}

\subsubsection{Workflow and Step History}
\begin{itemize}
    \item 4.6.1 The system shall display task and step history for each workflow.
    \item 4.6.2 History shall be shown in chronological order.
    \item 4.6.3 The system shall allow viewing history at both workflow and step levels in a
    readable format.
\end{itemize}

\subsection{Logical Database Requirements}

\subsubsection{Core Entities}
\begin{itemize}
    \item \textbf{User:} UserID, Username, Email, Role, Status
    \item \textbf{Workflow:} WorkflowID, Name, Description, Status, CreatedBy, CreatedDate
    \item \textbf{WorkflowStep:} StepID, WorkflowID, StepNumber, StepName, Description
    \item \textbf{Task:} TaskID, WorkflowID, StepID, AssignedTo, Status, StartDate,
    CompletionDate
\end{itemize}

\subsubsection{Constraints}
\begin{itemize}
    \item Each workflow must have at least one step.
    \item StepNumber must be unique within a workflow.
    \item Workflow names must be unique.
    \item Deleting a workflow must not delete dependent tasks.
\end{itemize}

\subsubsection{Data Retention}
\begin{itemize}
    \item Completed workflows and task records must be retained for three years.
    \item Archived data will be read-only.
\end{itemize}

\subsection{Design Constraints}
\begin{itemize}
    \item Built using ASP.NET MVC, C\#, Entity Framework.
    \item Must run under IIS on Windows Server.
    \item Supports SQL Server 2019+ and Oracle 19c+.
    \item HTTPS/TLS 1.2+ required.
    \item Up to 100 concurrent users; page load $\leq$ 3 seconds.
    \item Developed in Visual Studio 2022+.
    \item Version control through GitHub Enterprise.
    \item Only approved NuGet packages permitted.
\end{itemize}

% ─────────────────────────────────────────────
\section{Other Nonfunctional Requirements}
% ─────────────────────────────────────────────

\subsection{Performance Requirements}
\begin{itemize}
    \item Average system response time should not exceed 2--3 seconds for any database query
    under normal conditions.
    \item The system shall handle concurrent requests from at least 100 users.
    \item Workflow search queries should return results within 1 second for up to 10,000
    records.
\end{itemize}

\subsection{Safety Requirements}
\begin{itemize}
    \item The system shall prevent unauthorized deletion or modification of workflow data.
    \item The application shall include automated daily backups of all critical data.
    \item In case of failure, the system shall restore from the last known good backup.
\end{itemize}

\subsection{Security Requirements}
\begin{itemize}
    \item Authorization is based on a URL query string flag.
    \item Role-based access control (RBAC) must be strictly enforced.
    \item Session timeout will occur after 15 minutes of inactivity.
    \item All sensitive fields (e.g., user credentials, role mappings) must be encrypted in
    transit and at rest.
\end{itemize}

\subsection{Software Quality Attributes}
\begin{itemize}
    \item \textbf{Reliability:} The system must maintain 99.5\% uptime during business hours.
    \item \textbf{Usability:} Interfaces must follow accessibility (Section 508) guidelines.
    \item \textbf{Maintainability:} Source code should follow SOLID design principles and
    include XML documentation.
    \item \textbf{Scalability:} The system must allow future integration with third-party APIs.
    \item \textbf{Portability:} The application must support deployment in both SQL Server and
    Oracle environments.
\end{itemize}

\subsection{Business Rules}
\begin{itemize}
    \item Only administrators can create or modify workflows.
    \item Workflows cannot be deleted once created; they may only be archived.
    \item Tasks cannot be reassigned once marked as ``Completed.''
    \item Workflow names must be unique and descriptive to prevent duplication.
    \item Each workflow must contain at least one step before activation.
\end{itemize}

% ─────────────────────────────────────────────
\section{Legal and Ethical Considerations}
% ─────────────────────────────────────────────

\subsection{Scope and Sensitivity}
The Workflow Queue (WFQ) system may process sensitive organizational data and personally
identifiable information (PII) associated with authenticated users. The design therefore
prioritizes confidentiality, integrity, availability, and accountability in both data handling
and system behavior.

\subsection{Access Control and Authorization}
\begin{itemize}
    \item All users authenticate through the organization's identity provider.
    \item Role-based access control (RBAC) enforces least-privilege access; only users with
    the Admin role may access administrative functions.
    \item Administrative actions are logged to support accountability and post-event review.
\end{itemize}

\subsection{Transmission and Session Security}
\begin{itemize}
    \item All browser--server and service-to-service communications shall use HTTPS/TLS.
    \item Sessions shall follow organization policy for timeout and re-authentication (e.g.,
    15 minutes of inactivity).
    \item Sensitive values shall never be transmitted or stored in plaintext.
    \item Tokens, cookies, and session identifiers shall be marked HttpOnly and Secure and
    protected against CSRF.
\end{itemize}

\subsection{Data Integrity, Auditability, and Non-Repudiation}
\begin{itemize}
    \item Workflow and task history shall be preserved and viewable to authorized users.
    \item Deleting workflows in the UI shall not be supported; archival is the approved
    lifecycle state.
    \item Database operations shall avoid cascade deletes that would remove dependent records.
    \item All administrative actions shall be logged with timestamp, actor, and outcome.
\end{itemize}

\subsection{Retention and Archival}
\begin{itemize}
    \item Completed workflows and task records shall be retained per organizational retention
    policy and then archived in read-only form.
    \item Automated daily backups and documented restore procedures shall support continuity
    and evidence preservation.
    \item Access to archives shall be restricted and monitored.
\end{itemize}

\subsection{Compliance Alignment}
\begin{itemize}
    \item The system shall align with internal cybersecurity, authentication, and
    records-management policies.
    \item If scope later includes regulated data types (e.g., PHI), controls shall be extended
    to meet the applicable regulatory framework without weakening existing safeguards.
    \item Periodic reviews shall verify continued alignment with policy and regulatory changes.
\end{itemize}

% ─────────────────────────────────────────────
\appendix
\section{Glossary}
% ─────────────────────────────────────────────

\begin{center}
\begin{tabular}{ll}
\toprule
Term & Definition \\
\midrule
WFQ   & Workflow Queue System \\
RBAC  & Role-Based Access Control \\
ORM   & Object-Relational Mapping \\
HTTPS & Hypertext Transfer Protocol Secure \\
AD    & Active Directory \\
MVC   & Model-View-Controller \\
CRUD  & Create, Read, Update, Delete \\
\bottomrule
\end{tabular}
\end{center}

\section{Analysis Models}
\begin{itemize}
    \item \textbf{Use Case Diagram:} Displays relationships between users (Admin, Manager,
    User) and system functions (Create Workflow, Approve Task, Track Progress).
    \item \textbf{ER Diagram:} Shows logical relationships between User, Workflow,
    WorkflowStep, and Task entities.
    \item \textbf{Sequence Diagram:} Illustrates user login and workflow approval sequence.
\end{itemize}

\section{To Be Determined List}
\begin{center}
\begin{tabular}{lll}
\toprule
Item & Description & Target Decision Date \\
\midrule
Workflow Name Max Length        & Maximum character limit for Name field              & TBD by Dev Team \\
StepNumber Auto-Increment Rule  & Define logic for assigning new step numbers         & TBD by Design Lead \\
User Role Mapping Details       & Final mapping between AD roles and WFQ roles        & TBD by QTC Security Team \\
Reporting Module                & Inclusion of dashboard analytics                    & Version 2.0 \\
Notification Integration        & Email and Teams alerts for approvals                & Version 2.0 \\
\bottomrule
\end{tabular}
\end{center}

\end{document}
